\subsubsection{Yêu cầu của sản phẩm được triển khai}
\label{sec:t3-requirements}

Mục này chốt phạm vi phần mềm mà Task 3 triển khai. Yêu cầu được rút ra từ phạm vi MVP ở Task 1.4, Client Interaction Log và các quyết định làm rõ trong quá trình thực hiện. Mỗi yêu cầu đã được \textbf{đối chiếu lại với mã nguồn bàn giao} (backend Spring Boot, frontend Angular và tài liệu \texttt{docs/api}) để bảo đảm báo cáo mô tả đúng hệ thống đang chạy, không mô tả theo bản thiết kế ban đầu. Mã FR và NFR được dùng lại trong ma trận truy vết và bằng chứng kiểm thử ở Mục \ref{sec:t3-quality}.

\subsubsubsection{Yêu cầu chức năng (Functional Requirements)}

\begingroup\small
\renewcommand{\arraystretch}{1.15}
\begin{longtable}{@{}p{.08\linewidth}p{.52\linewidth}p{.13\linewidth}p{.17\linewidth}@{}}
\caption{Danh mục yêu cầu chức năng}\label{tab:t3-fr}\\
\toprule
\textbf{Mã} & \textbf{Yêu cầu} & \textbf{Use case} & \textbf{Increment / HĐ}\\\midrule\endfirsthead
\toprule
\textbf{Mã} & \textbf{Yêu cầu} & \textbf{Use case} & \textbf{Increment / HĐ}\\\midrule\endhead
\midrule\multicolumn{4}{r}{\textit{Còn tiếp}}\\\endfoot
\bottomrule\endlastfoot
\multicolumn{4}{@{}l}{\textit{Nhóm 1 -- Tài khoản và phân quyền}}\\
FR-01 & Sinh viên tự tạo tài khoản (họ tên, MSSV, email, khoa, năm học, mật khẩu); vai trò luôn là Sinh viên, MSSV và email không trùng. & UC-01 & I1 / A05.1\\
FR-02 & Đăng nhập bằng email và mật khẩu, đăng xuất, khôi phục phiên khi tải lại trang. & UC-02, UC-03 & I1 / A05.1\\
FR-03 & Đổi mật khẩu; tài khoản dùng mật khẩu tạm bắt buộc đổi trước khi dùng các chức năng khác. & UC-25 & I1 / A05.1\\
FR-04 & Ba vai trò cố định (Sinh viên, Người tổ chức, Quản trị viên); mọi API kiểm tra quyền ở backend; Quản trị viên không mặc nhiên có quyền Người tổ chức. & Tất cả & I1 / A05.1\\
FR-05 & Quản trị viên tìm kiếm tài khoản theo MSSV/email/họ tên, lọc theo vai trò và trạng thái. & UC-21 & I1 / A05.1\\
FR-06 & Quản trị viên cấp tài khoản Người tổ chức với mật khẩu tạm và lý do. & UC-22 & I1 / A05.1\\
FR-07 & Quản trị viên khóa/mở khóa tài khoản có lý do; không tự khóa; luôn còn ít nhất một Quản trị viên đang hoạt động. & UC-23 & I1 / A05.1\\
\multicolumn{4}{@{}l}{\textit{Nhóm 2 -- Sự kiện}}\\
FR-08 & Người tổ chức tạo sự kiện nháp với đủ thông tin: tên, mô tả, địa điểm, đơn vị, 4 mốc thời gian, sức chứa nguyên dương, giới hạn khoa/năm (chọn nhiều). & UC-12 & I1 / A05.2\\
FR-09 & Người tổ chức sửa sự kiện của mình; kiểm tra xung đột phiên bản; chặn thay đổi làm sai đăng ký hiện có. & UC-13 & I1, I2 / A05.2, A07.1\\
FR-10 & Người tổ chức công bố sự kiện nháp hợp lệ; chỉ sự kiện đã công bố mới hiển thị với người dùng khác. & UC-14 & I1 / A05.2\\
FR-11 & Người dùng đã đăng nhập tìm sự kiện đã công bố theo tên (phân trang) và xem chi tiết. & UC-04, UC-05 & I1 / A05.2\\
\multicolumn{4}{@{}l}{\textit{Nhóm 3 -- Đăng ký, danh sách chờ và email}}\\
FR-12 & Sinh viên đủ điều kiện đăng ký trong khung giờ đăng ký; chỉ người \texttt{REGISTERED} chiếm chỗ; không đăng ký trùng. & UC-06 & I2 / A07.1\\
FR-13 & Khi hết chỗ, sinh viên được đưa vào danh sách chờ (\texttt{WAITLISTED}) và biết vị trí của mình. & UC-07 & I2 / A07.2\\
FR-14 & Sinh viên hủy đăng ký hoặc rời danh sách chờ trước giờ bắt đầu; được đăng ký lại trong khung giờ đăng ký. & UC-08 & I2 / A07.1\\
FR-15 & Hệ thống tự đôn người chờ hợp lệ đầu tiên (FIFO) khi có chỗ trống, kể cả sau giờ đóng đăng ký, cho đến giờ bắt đầu. & UC-09 & I2 / A07.2\\
FR-16 & Gửi email khi: đăng ký có chỗ, được đôn từ danh sách chờ, sự kiện đổi địa điểm hoặc thời gian. & UC-06, UC-09, UC-13 & I2 / A07.3\\
FR-17 & Sinh viên xem danh sách sự kiện của mình theo trạng thái và kết quả điểm danh của mình. & UC-10 & I2, I3 / A07.1\\
FR-18 & Người tổ chức xem danh sách đăng ký và danh sách chờ của sự kiện mình phụ trách. & UC-15 & I2 / A07.1\\
\multicolumn{4}{@{}l}{\textit{Nhóm 4 -- Điểm danh, báo cáo và xuất dữ liệu}}\\
FR-19 & Người tổ chức mở/đóng phiên QR có thời hạn 1--30 phút, trong khoảng từ 30 phút trước giờ bắt đầu đến giờ kết thúc. & UC-16 & I3 / A09.1\\
FR-20 & Sinh viên \texttt{REGISTERED} quét QR, xem trước thông tin rồi bấm xác nhận để ghi nhận có mặt; mỗi người một lượt. & UC-11 & I3 / A09.1\\
FR-21 & Người tổ chức ghi nhận hoặc sửa điểm danh thủ công, bắt buộc có lý do; mọi thay đổi được lưu lịch sử. & UC-17, UC-18 & I3 / A09.2\\
FR-22 & Báo cáo theo sự kiện: số đăng ký, chờ, hủy, có mặt, vắng và tỷ lệ tham dự; số vắng chỉ chốt sau giờ kết thúc. & UC-19 & I3 / A10.1\\
FR-23 & Xuất báo cáo và danh sách chi tiết ra tệp CSV đọc được tiếng Việt. & UC-20 & I3 / A10.2\\
\end{longtable}
\endgroup

\subsubsubsection{Yêu cầu phi chức năng (Non-functional Requirements)}

\begingroup\small
\renewcommand{\arraystretch}{1.15}
\begin{longtable}{@{}p{.08\linewidth}p{.14\linewidth}p{.44\linewidth}p{.25\linewidth}@{}}
\caption{Danh mục yêu cầu phi chức năng}\label{tab:t3-nfr}\\
\toprule
\textbf{Mã} & \textbf{Thuộc tính} & \textbf{Yêu cầu đo được} & \textbf{Cách hiện thực / kiểm chứng}\\\midrule\endfirsthead
\toprule
\textbf{Mã} & \textbf{Thuộc tính} & \textbf{Yêu cầu đo được} & \textbf{Cách hiện thực / kiểm chứng}\\\midrule\endhead
\midrule\multicolumn{4}{r}{\textit{Còn tiếp}}\\\endfoot
\bottomrule\endlastfoot
NFR-01 & Hiệu năng & Hiển thị báo cáo $\le$ 5 giây và tạo tệp CSV $\le$ 10 giây với 500 bản ghi đăng ký (chỉ số M2, Task 1.4). & Truy vấn một snapshot; đo trên bộ 2.000 đăng ký (Mục \ref{sec:t3-quality}).\\
NFR-02 & Toàn vẹn dữ liệu & Không vượt sức chứa, không đăng ký trùng, không đôn sai thứ tự khi có nhiều yêu cầu đồng thời; lỗi giữa chừng không để lại dữ liệu dở dang. & Khóa sự kiện trong transaction, \texttt{UNIQUE(event\_id, student\_id)}, rollback toàn bộ; test tranh chấp 6 người/1 chỗ.\\
NFR-03 & Chống ghi đè & Hai người cùng sửa một bản ghi thì người lưu sau nhận lỗi xung đột thay vì ghi đè âm thầm. & Optimistic locking bằng \texttt{version}, trả 409 \texttt{VERSION\_CONFLICT}.\\
NFR-04 & Bảo mật xác thực & Mật khẩu băm BCrypt (cost 12); phiên lưu trong cookie HttpOnly, SameSite=Lax, hết hạn 8 giờ; không lưu token ở localStorage. & \texttt{auth\_sessions} chỉ lưu SHA-256 của token phiên.\\
NFR-05 & Bảo mật ứng dụng & Mọi request ghi có CSRF token; giới hạn tần suất đăng nhập/đăng ký (429); truy cập tài nguyên ngoài quyền trả 404 để không lộ sự tồn tại. & Test RBAC/CSRF ở I1--I3; rate limit kiểm thử bằng unit test và HTTP.\\
NFR-06 & An toàn QR & Mã QR ngẫu nhiên 32 byte, chỉ hợp lệ trong phiên còn hạn; DB không lưu token thô. & Lưu SHA-256 và AES-256-GCM; guard ở database kiểm tra hạn ngay lúc ghi.\\
NFR-07 & Tin cậy email & Lỗi email không làm hỏng thao tác nghiệp vụ; tự thử lại tối đa 5 lần với lỗi tạm thời; không gửi thư cho thao tác đã rollback. & Transactional outbox + Brevo HTTP API.\\
NFR-08 & Riêng tư dữ liệu & Dùng dữ liệu sinh viên giả lập; sinh viên chỉ xem dữ liệu của chính mình; Người tổ chức chỉ xem sự kiện mình phụ trách. & Định danh lấy từ phiên, không nhận \texttt{studentId} từ client.\\
NFR-09 & Khả dụng & Giao diện tiếng Việt, hiển thị giờ Việt Nam (UTC+07:00); dùng được trên điện thoại để quét QR. & Angular responsive; camera cần HTTPS.\\
NFR-10 & Bảo trì & Backend phân tầng Controller--Service--Repository; schema được quản lý bằng migration; có pipeline CI. & Flyway V1/V2; \texttt{.github/workflows/validation.yml}.\\
NFR-11 & Giới hạn tài nguyên & Một lần xuất CSV tối đa 10.000 dòng; vượt quá trả 413, không cắt dữ liệu âm thầm. & Kiểm thử 10.001 dòng bị chặn.\\
NFR-12 & Triển khai & Chạy được trên môi trường cloud qua HTTPS và có hướng dẫn chạy lại ở local bằng Docker Compose. & Vercel + Render + Neon; \texttt{README.md}, \texttt{deploy/}.\\
\end{longtable}
\endgroup
